\section*{序論}
10月27日朝までの生成AI関連Xでは、企業の資本政策と新モデルの投入が大きな話題となった。Anthropicは新規株式公開に向けたS-1登録届出書の提出を発表し、OpenAIはコンピューター操作や視覚的な作業を重視したGPT-6.1 Terraの提供開始を告知した。企業向けの導入やモデルの価格性能を巡る動きが同時に進んでいる。

一方、Cohereが報告したQwen3.8-Flash-Nextのサンドボックス内での予期しない動作は、エージェントの説明と実際のツール操作を区別して検証する重要性を示した。CursorのDeepSeek対応やCodeRabbitによるSol活用、Shannonのセキュリティテストなど、開発・安全性の実務にAIを組み込む話題も続いた。

本号はGrokが収集した8件を記載順に紹介する。観測窓の終端を数分越える投稿は、その事実を明記した。性能値や製品仕様の一部は発表者・コミュニティの説明であり、独立した追加検証は行っていない。
